% 来源：成文/A\_04\_问题重述.md
\section{问题重述}
中药材的干燥是决定药材品质的关键工序，热风烘干是产地常用的方式。就数学形式而言，题目要求求解的是一类以药材内部温度场与干基含水率场为未知函数的\textbf{传热传质耦合非稳态问题}：两场沿半径方向变化，并随烘房环境随时间演化；当物性随状态变化或几何域发生收缩时，控制方程表现为非线性。题目要求针对圆柱形药材建立描述烘干过程的数学模型，并依次解决以下四个问题。

\textbf{问题一：} 给定药材初始温度 28 ℃、初始干基含水率 2.55 kg/kg，以及烘房温度与水分浓度的实测时序，要求刻画预热平衡阶段 0–1800 s 内药材内部温度场与水分浓度场的分布规律，并按题面规定的时刻与位置输出结果。

\textbf{问题二：} 物性改为随含水率与温度变化（相关经验式按附录 3 取用），要求刻画整个烘干过程 0–3 h 内两场的演化规律。与问题一的实质差别在于物性的状态依赖使两场由解耦转为耦合，输出要求与问题一相同。

\textbf{问题三：} 在要求药材内部各处干基含水率均低于 0.15 kg/kg 的前提下，反演满足该要求的烘干时间（单位：h），并给出指定时刻与位置的水分浓度。该问的实质是带终端判据的反问题。

\textbf{问题四：} 药材在烘干中会因失水而收缩，要求依据附件 2 给出的半径时序（相关经验式按附录 4 取用）确定其烘干时长，并在输出中给出包括药材表面在内的水分浓度。与问题三的实质差别在于计算域随时间收缩。

\subsection{建模对象与给定数据}
本节把建模对象与给定数据一并交代（见表~\ref{tab:1-1}），后文各章不再重复。



\settabw{4}
\begin{longtable}{>{\raggedleft\arraybackslash}m{\dimexpr 0.196\tblw\relax}>{\raggedright\arraybackslash}m{\dimexpr 0.804\tblw\relax}}
  \caption{建模对象与给定数据}\label{tab:1-1}\\
  \toprule
    项 & 内容 \\
  \midrule
  \endfirsthead
  \multicolumn{2}{r}{\footnotesize（续）}\\
  \toprule
    项 & 内容 \\
  \midrule
  \endhead
  \bottomrule
  \endfoot
  \bottomrule
  \endlastfoot
  \textbf{几何} & 圆柱形药材，长度 25 cm、初始半径 2 cm；长径比足够大，只考察半径方向的一维轴对称问题（见假设 1） \\
  \textbf{初值} & 温度 28 ℃、干基含水率 2.55 kg/kg，全域均匀（题面仅给单一初值） \\
  \textbf{物性} & 题面\textbf{分套给出}：问题一按附录 2（常物性），问题二、三按附录 3，问题四按附录 4；表面换热系数 25 W/(m\textsuperscript{2}·K) 与对流传质系数 $8\times10^{-7}$ m/s 由附录 2 给出并全程沿用 \\
  \textbf{环境数据} & 附件 1：0–14400 s 的烘房温度与水分浓度实测时序（241 点，间隔 60 s）；其后以恒定值外推 \\
  \textbf{收缩数据} & 附件 2：半径时序（145 点，间隔 1800 s），由 2.000 cm 收缩至 1.198 cm（仅问题四使用） \\
  \textbf{达标判据} & 全域干基含水率\textbf{严格小于} 0.15 kg/kg \\
  \textbf{输出契约} & 题面规定的六张结果表按题面行列给出（温度与水分浓度各三张；正式呈现见第五章相应各节），其中前两表以 s 计时、后四表以 h 计时，末表末列为药材表面；完整结果另存为四个结果文件，A 列一律为 s、结果保留四位小数 \\
\end{longtable}

题面所述“低于 0.15 kg/kg”按\textbf{严格小于}执行，即在时间步级判定全域最大干基含水率小于该阈值，判定与终止方式见第五章问题三一节。
